\documentclass[11pt]{article}
\usepackage[T1]{fontenc}
\usepackage{lmodern,amsmath,amssymb,amsthm}
\usepackage[margin=1in]{geometry}
\usepackage[colorlinks=true,linkcolor=blue,citecolor=blue,urlcolor=blue]{hyperref}
\usepackage{microtype}
\hypersetup{pdftitle={A priori estimates across the subcritical Lane--Emden hyperbola},pdfauthor={Alec Kriebel},pdfsubject={Consequences of the OpenAI subcritical Liouville theorem}}
\newcommand{\R}{\mathbb R}
\newcommand{\dist}{\operatorname{dist}}
\newcommand{\norm}[1]{\left\lVert#1\right\rVert}
\newtheorem{theorem}{Theorem}
\newtheorem{lemma}[theorem]{Lemma}
\newtheorem{remark}[theorem]{Remark}
\title{A priori estimates across the subcritical\texorpdfstring{\\}{ }Lane--Emden hyperbola}
\author{Alec Kriebel\\\small\href{https://orcid.org/0009-0001-9320-500X}{ORCID: 0009-0001-9320-500X}}
\date{October 6, 2026}
\begin{document}
\maketitle
\begin{abstract}
We record the consequences of the subcritical entire Liouville theorem in
OpenAI's 2026 H\'enon--Lane--Emden preprint for the unweighted system
$-\Delta u=v^p$, $-\Delta v=u^q$. For $n\ge3$, $p,q>1$, and
$1/(p+1)+1/(q+1)>(n-2)/n$, the established reductions of
Pol\'a\v{c}ik--Quittner--Souplet give universal pointwise estimates on
arbitrary proper domains and zero-Dirichlet half-space nonexistence.
We detail the resulting gradient and exterior estimates and the
boundary blow-up argument giving uniform bounds and compactness on
fixed bounded $C^2$ domains, with exact function spaces.
The full subcritical scope is an immediate consequence of the cited
Liouville input; no independent solution of that conjecture or new
Liouville-to-estimate machinery is claimed.
\end{abstract}

\section{Scope, provenance, and statement}
Let
\begin{equation}\label{range}
 n\ge3,\qquad p,q>1,\qquad
 \frac1{p+1}+\frac1{q+1}>\frac{n-2}{n},
 \qquad
 \alpha=\frac{2(p+1)}{pq-1},\quad
 \beta=\frac{2(q+1)}{pq-1}.
\end{equation}
Throughout, a domain is a nonempty connected open subset of $\R^n$.
A nonnegative classical solution means $u,v\in C^2(\Omega)$,
$u,v\ge0$, satisfying
\begin{equation}\label{system}
 -\Delta u=v^p,\qquad -\Delta v=u^q
 \quad\text{pointwise in }\Omega.
\end{equation}
Zero-Dirichlet solutions additionally belong to $C(\overline\Omega)$
and vanish on $\partial\Omega$. These statements do not assert
regularity for arbitrary weak or distributional solutions.

The entire Liouville theorem used here is the unweighted corollary of
OpenAI \cite{OAI}, at the source version specified below.
The passage from bounded entire nonexistence to the universal domain
and half-space conclusions is already Theorems 4.2 and 4.3 of
Pol\'a\v{c}ik--Quittner--Souplet \cite{PQS} (PQS).
The fixed-domain Dirichlet reduction is also established; see
Quittner--Souplet \cite[Theorem 4.1(ii) and Remark 4.1(ii)]{QS},
whose remark credits the pure-power case to PQS.
We include its blow-up argument to distinguish a boundary estimate
from a uniform Dirichlet bound and to specify compactness.

Thus the contribution of this consequence note is the explicit full-range
formulation and checked compatibility of the input, boundary argument,
and function spaces. Dimensions three and four, and portions of higher
dimensions, already had the needed Liouville results; see
\cite{PQS,Souplet,LiLiWei}. The newly available scope is the remaining
portions of \eqref{range}, inherited from \cite{OAI}, rather than a
new estimate formula. We do not claim priority for the transfer method
or to be first to observe its implication. The source preprint carries
September 24, 2026 as its manuscript date; the verified public release
is October 6, 2026 \cite{release}. Our priority audit found no earlier
explicit full-range application package in that release, but absence
from a search does not certify priority.

\begin{theorem}\label{main}
Assume \eqref{range}. There is $C=C(n,p,q)>0$ such that every
nonnegative classical solution of \eqref{system} on every proper
domain $\Omega\ne\R^n$ satisfies, with $d(x)=\dist(x,\partial\Omega)$,
\begin{equation}\label{distance}
 \begin{aligned}
 u(x)&\le C d(x)^{-\alpha},&
 |\nabla u(x)|&\le C d(x)^{-\alpha-1},\\
 v(x)&\le C d(x)^{-\beta},&
 |\nabla v(x)|&\le C d(x)^{-\beta-1}.
 \end{aligned}
\end{equation}
No boundary regularity or boundary condition is required.
If $\Omega\supset\{x:|x|>R\}$ for some $R>0$, then for
$|x|\ge2R$ the same estimates hold with $d(x)$ replaced by $|x|$
(after changing $C$).

Every nonnegative zero-Dirichlet classical solution on
$H=\{x_n>0\}$ in $C^2(H)\cap C(\overline H)$ is zero, with no
boundedness or growth assumption.

For each fixed bounded $C^2$ domain $\Omega$, its nonnegative
zero-Dirichlet classical solutions satisfy
\begin{equation}\label{dirichlet}
 \norm{u}_{L^\infty(\Omega)}+\norm{v}_{L^\infty(\Omega)}
 \le C(\Omega,n,p,q).
\end{equation}
For every $1<r<\infty$ and $0<\gamma<1$, this solution set is bounded
and sequentially compact in
\[
 [W^{2,r}(\Omega)\cap W^{1,r}_0(\Omega)]^2
 \quad\text{and}\quad [C^{1,\gamma}(\overline\Omega)]^2,
\]
respectively, with the indicated norm topologies.
If $\partial\Omega$ is $C^{2,\theta}$, $0<\theta<1$, there are uniform
$C^{2,\theta}(\overline\Omega)$ bounds, and sequential compactness
in $[C^{2,\eta}(\overline\Omega)]^2$ for every $0<\eta<\theta$.
Constants for these norm bounds may depend on the fixed domain and the
indicated regularity exponent. No $C^2(\overline\Omega)$ conclusion is
asserted from $C^2$ boundary regularity alone.
\end{theorem}

\section{Entire input and interior normalization}
\begin{lemma}\label{entire}
Under \eqref{range}, the only bounded nonnegative classical solution
of \eqref{system} on $\R^n$ is $(0,0)$.
\end{lemma}
\begin{proof}
The unweighted corollary of \cite{OAI} excludes strictly positive
$C^2(\R^n)$ solution pairs, even without boundedness. If either
nonnegative component vanishes at a point, the strong maximum principle
for its superharmonic equation makes it identically zero. Its equation
then forces the other component to be zero. Otherwise both components
are strictly positive, and the cited theorem applies.
\end{proof}

The source input is substantial, not a hypothesis hidden in notation.
Our dependency audit inspected its Newton-potential reductions,
localized virial identity, interval comparison, and energy-relative
layer-loss estimate before the scaling contradiction. The accompanying
package gives the exact source hashes and a reconstruction of the
unweighted argument. The actual Lean source declaration is separately
identified in the audit; its formal-verification status is reported
there. A complete pinned Lean kernel build was not reproduced, and
this follow-on theorem is not claimed to be formalized. The mathematical
dependency acceptance here rests on the analytic proof audit.

For clarity, we recall the normalization underlying PQS Theorem 4.3.
The scaling identities are
\begin{equation}\label{scaling}
 p\beta=\alpha+2,\quad q\alpha=\beta+2,
 \qquad U(y)=\lambda^\alpha u(x+\lambda y),\quad
 V(y)=\lambda^\beta v(x+\lambda y).
\end{equation}
The pair $(U,V)$ satisfies the same equations. Put
$M=u^{1/\alpha}+v^{1/\beta}$. For a nonzero solution, the maximum
principle ensures $M>0$. If no universal bound for $Md$ exists,
choose solutions with $M(y_k)d(y_k)>2k$.
The doubling lemma \cite[Lemma 5.1]{PQS} selects $x_k$ with
\[
 M(x_k)\ge M(y_k),\quad M(x_k)d(x_k)>2k,\quad
 M(z)\le2M(x_k)\quad (|z-x_k|\le k/M(x_k)).
\]
One can see the selection directly: whenever the last property fails,
move by at most $k/M$ to a point with more than twice that value.
The sum of movements is at most $2k/M(y_k)$, strictly less than
$d(y_k)$; if the process never stops, local boundedness at its
interior limit contradicts $M\to\infty$.

With $\lambda_k=M(x_k)^{-1}$, the rescaled domains contain $B_k$,
$U_k^{1/\alpha}(0)+V_k^{1/\beta}(0)=1$, and
$U_k^{1/\alpha}+V_k^{1/\beta}\le2$ on $B_k$.
Interior elliptic estimates first give local $W^{2,r}$ bounds,
then $C^{1,\gamma}$ bounds for suitable $r>n$. Since $p,q>1$, the
powers are Lipschitz on bounded nonnegative ranges. Interior Schauder
estimates therefore give uniform $C^{2,\gamma}$ bounds on smaller
balls. A diagonal subsequence converges in $C^2_{\mathrm{loc}}(\R^n)$
to a bounded nonnegative pair satisfying the equations and the same
nonzero normalization. This contradicts Lemma \ref{entire}.
Consequently $Md\le C$, which proves the pointwise part of
\eqref{distance} with constants depending only on $n,p,q$.

\section{Gradients, exterior decay, and the half-space}
Fix $x\in\Omega$ and $\rho=d(x)/2$. On $B_\rho(x)$ the distance
is at least $d(x)/2$. Applying the pointwise bounds to the scaling
\eqref{scaling} with $\lambda=\rho$ gives uniformly bounded
$U,V$ and right sides on $B_1$. The interior gradient estimate for
Poisson's equation yields
\[
 |\nabla U(0)|\le C(\norm{U}_{L^\infty(B_1)}+
                           \norm{V^p}_{L^\infty(B_1)})\le C,
\]
and similarly for $V$. Rescaling proves the gradient part of
\eqref{distance}. For the exterior conclusion, every boundary point
lies in $\overline B_R$, so $d(x)\ge |x|-R\ge |x|/2$ when
$|x|\ge2R$.

PQS Theorem 4.2 assumes exactly the bounded nonnegative entire
nonexistence of Lemma \ref{entire}, with $p,q>1$, and excludes every
nontrivial nonnegative $C^2(H)\cap C(\overline H)$ zero-Dirichlet
solution, bounded or unbounded. This proves the asserted half-space
conclusion. Its full conclusion is needed here; an arbitrary-domain
distance estimate by itself diverges at the boundary and would not
prove \eqref{dirichlet}.

\section{Fixed-domain boundary blow-up and compactness}
We give the remaining argument explicitly. On a bounded $C^2$ domain,
a classical zero-Dirichlet pair continuous on the closure already
belongs to $W^{2,r}(\Omega)\cap W^{1,r}_0(\Omega)$ for every
$1<r<\infty$: its right sides are bounded; solve the corresponding
linear Dirichlet problems, and compare by the maximum principle.
Indeed the difference is harmonic, continuous on the closure, and zero
on the boundary. The global Poisson estimates on a $C^{1,1}$ domain
(which a bounded $C^2$ domain is) give the stated regularity. We use
also their local boundary versions and Sobolev embeddings; standard
forms are in \cite{GT}.

Suppose \eqref{dirichlet} fails. For a sequence of solutions set
\[
 A_k=\max_{\overline\Omega}
       (u_k^{1/\alpha}+v_k^{1/\beta})\longrightarrow\infty,
 \quad M_k(x_k)=A_k,\quad \lambda_k=A_k^{-1}.
\]
The maximizing points are interior. Scaling about $x_k$ gives
$0\le U_k,V_k\le1$ and
$U_k^{1/\alpha}(0)+V_k^{1/\beta}(0)=1$.
If $d(x_k)/\lambda_k\to\infty$ along a subsequence, the interior
argument just given produces a nontrivial bounded entire solution,
a contradiction.

Otherwise pass to a subsequence with
$t_k=d(x_k)/\lambda_k\to t\ge0$. Then $d(x_k)\to0$.
Choose a nearest boundary point $b_k$ and rotate its inward normal
into the $y_n$ direction. Uniform $C^2$ boundary charts exist by
compactness of $\partial\Omega$. Locally they have graphs
$x_n=\varphi_k(x')$, with $\varphi_k(0)=D\varphi_k(0)=0$ and
uniformly bounded second derivatives. After scaling at $b_k$, the
graphs are
\[
 y_n=\psi_k(y')=\lambda_k^{-1}\varphi_k(\lambda_k y'),
\]
so $\psi_k\to0$ in $C^2$ on every fixed compact set; the point $x_k$
becomes $(0,t_k)$. Pull back by
$(y',y_n)\mapsto(y',y_n+\psi_k(y'))$.
The resulting uniformly elliptic coefficients tend to those of the
Laplacian, and the lower order coefficients tend to zero.
The charts have uniform $C^{1,1}$ control, the solutions and right
sides are uniformly bounded, and the flat-boundary trace is zero.
Local boundary $W^{2,r}$ estimates therefore give uniform
$C^{1,\gamma}$ bounds on smaller compact half-balls for any
$0<\gamma<1$, choosing $r$ sufficiently large.

A diagonal subsequence converges in $C^1$ on compact subsets of
$\overline H$ and in $C^2$ on compact subsets of $H$ to a bounded
nonnegative solution $(U,V)$, continuous and zero on $\partial H$.
The transformed equations converge in distributions to the Laplacian
equations; interior regularity upgrades the limit to $C^2(H)$.
Normalization passes at $(0,t)$. If $t=0$, continuity and the zero
trace contradict normalization. If $t>0$, the limiting pair is
nontrivial, contradicting the half-space theorem. Both cases exclude
the assumed blow-up, proving \eqref{dirichlet}.

The uniform $L^\infty$ bound now gives uniform global $W^{2,r}$ bounds
for every finite $r>1$. Compact Sobolev embeddings, using a larger
$r$ if necessary, give a subsequence converging in
$C^{1,\gamma}(\overline\Omega)$ for any fixed $\gamma<1$.
The limiting pair retains nonnegativity and zero trace, and its equations
hold in distributions and then classically in the interior. The right
sides converge uniformly. For each finite $r>1$, the global linear
Dirichlet estimate applied to differences gives
\[
 \norm{u_k-u}_{W^{2,r}}\le C_r\norm{v_k^p-v^p}_{L^r}\to0,
 \qquad
 \norm{v_k-v}_{W^{2,r}}\le C_r\norm{u_k^q-u^q}_{L^r}\to0.
\]
Thus the compactness is strong, not just weak $W^{2,r}$ compactness.
On a $C^{2,\theta}$ domain, the already uniform Lipschitz bounds make
the right sides uniformly $C^{0,\theta}$; global Schauder estimates
give uniform $C^{2,\theta}$ bounds and compactness in $C^{2,\eta}$,
$\eta<\theta$. This completes Theorem \ref{main}.

\section{Limits and verification disclosure}
The strict hyperbola and $p,q>1$ are used as stated. No critical or
supercritical case, weighted boundary estimate, parabolic estimate,
or weak-solution regularity assertion is included. The conditional
asymptotically-power extension in PQS Theorem 7.3 is established work;
we do not present it as a new extension here. Spatially varying
coefficients are outside this note.

The upstream input is the public commit
\texttt{adc7f124} (full identity in the dependency ledger); source hashes,
analytic dependency checks, priority evidence, and complete-package
adversarial reviews accompany this note. Automated checks and reviews
are evidence about a proof, not a substitute for it. AI tools were used
extensively in research, drafting, and verification. This preprint has
not undergone conventional human peer review or refereeing at publication.
The author alone is listed; the upstream AI-produced theorem is
attributed to OpenAI using its manuscript-specific citation.

\begin{thebibliography}{9}
\bibitem{OAI} OpenAI, \emph{The Subcritical H\'enon--Lane--Emden
Conjecture}, OpenAI Math Release preprint (2026), manuscript dated
September 24, 2026, public release October 6, 2026.
\href{https://github.com/openai/math/tree/adc7f1241b42e322a6451854ab7e4b4c146bf78a/preprints/The-Subcritical-Henon-Lane-Emden-Conjecture-September-24-2026}{Pinned manuscript and supplied citation}.
\bibitem{PQS} P. Pol\'a\v{c}ik, P. Quittner, and Ph. Souplet,
\emph{Singularity and decay estimates in superlinear problems via
Liouville-type theorems. Part I: Elliptic equations and systems},
Duke Math. J. \textbf{139} (2007), 555--579.
\href{https://doi.org/10.1215/S0012-7094-07-13935-8}{doi:10.1215/S0012-7094-07-13935-8}.
The theorem statements were checked in the
\href{https://www-users.cse.umn.edu/~polacik/Publications/pqs1.pdf}{author-hosted preprint}.
\bibitem{QS} P. Quittner and Ph. Souplet,
\emph{Liouville theorems and universal estimates for superlinear
elliptic problems without scale invariance}, preprint,
\href{https://arxiv.org/abs/2407.04154v2}{arXiv:2407.04154v2} (2024),
Theorem 4.1(ii) and Remark 4.1(ii).
\bibitem{Souplet} Ph. Souplet, \emph{The proof of the Lane--Emden
conjecture in four space dimensions}, Adv. Math. \textbf{221} (2009),
1409--1427.
\bibitem{LiLiWei} K. Li, M. Li, and J. Wei,
\emph{On a new region for the Lane--Emden conjecture in higher dimensions},
\href{https://arxiv.org/abs/2510.06613v1}{arXiv:2510.06613v1} (2025).
\bibitem{GT} D. Gilbarg and N. S. Trudinger,
\emph{Elliptic Partial Differential Equations of Second Order},
Classics in Mathematics, Springer (2001), Chapters 6 and 9.
\bibitem{release} OpenAI, \emph{Sharing AI progress in mathematics},
October 6, 2026,
\url{https://openai.com/index/sharing-ai-progress-in-mathematics/}.
\end{thebibliography}
\end{document}
